\chapter{Testing and Debugging}
\label{ch:testing}

\epigraph{``Houston, we have a problem.''}{Jim Lovell, \emph{Apollo 13} (1995)}

What made Apollo 13 survivable was not heroism. It was that the ground had a
simulator, telemetry on everything, and the discipline to reproduce the failure
before attempting a fix. That is the whole content of this chapter, applied to
systems where one of the components is a language model and therefore cannot be
made deterministic.

The good news: almost none of MCP is the model. Almost all of it is a protocol,
and protocols are extremely testable.

\section{The testing pyramid, adapted}

\begin{table}[htbp]
\centering
\small
\begin{tabularx}{\textwidth}{@{}lXll@{}}
\toprule
\thead{Layer} & \thead{Tests} & \thead{Model} & \thead{Speed} \\
\midrule
Unit        & Handlers, schemas, serialisation & none & microseconds \\
Contract    & Conformance to \texttt{2026-07-28} & none & microseconds \\
Integration & Host to client to server flows & stubbed & milliseconds \\
Eval        & Does it actually work & real & minutes and money \\
\bottomrule
\end{tabularx}
\caption{Only the bottom layer needs a real model, and only the bottom layer is
slow. That is not a coincidence.}
\label{tab:test-pyramid}
\end{table}

Meridian's 191 tests run in 2.3 seconds and use no model at all. The eval suite is separate, slower, and gates deploys.

\keyidea{If your MCP tests are slow, you are testing the model when you meant to
test the protocol. Push everything you can down a layer.}

\section{Unit testing servers}

The trick that makes this pleasant: an in-process transport with the same
interface as the real ones.

\begin{py}
class InProcessTransport:
    """Talk to a `Server` object with no bytes on any wire.

    Serialisation is still performed, then immediately reversed. That is
    deliberate: skipping it would make the control group cheaper than any real
    transport could ever be, and the comparison would flatter the wrong things.
    """
\end{py}

Same \texttt{send(message) -> response} contract as stdio and HTTP, so the same
client code drives all three. Tests get microsecond round trips; benchmarks get
a control group. One class, two jobs.

\subsection{Assert on the wire}

The most important habit in this chapter.

\begin{py}
def test_server_identifies_itself_in_every_result(self):
    resp = self.server.handle(request("tools/list"))
    self.assertEqual(resp["result"]["_meta"][KEY_SERVER_INFO]["name"],
                     "test-server")
\end{py}

Not \texttt{server.info.name}. The bytes. Because the bytes are what another
implementation has to interoperate with, and an SDK convenience object can be
correct while the serialisation is broken.

\subsection{Test the round trip, not the encoder}

\begin{py}
def test_a_literal_sentinel_is_itself_encoded(self):
    literal = "=?base64?literal?="
    encoded = encode_header_value(literal)
    self.assertNotEqual(encoded, literal)
    self.assertEqual(decode_header_value(encoded), literal)
\end{py}

Two assertions: it \emph{did} something, and what it did is reversible. Testing
only the second would pass for an encoder that does nothing at all.

\section{Contract tests}

Conformance to \texttt{2026-07-28}, and this is the suite that will still be
useful when you rewrite the server.

Five families, and Meridian has all five.

\subsection{The envelope}

Does the server reject requests missing required \texttt{\_meta} fields? Does
it return an \texttt{Unsupported\-Protocol\-Version} error carrying a
\texttt{supported} list?

\begin{py}
def test_unsupported_version_lists_what_is_supported(self):
    resp = self.server.handle(request("tools/list", version="1999-01-01"))
    self.assertEqual(resp["error"]["code"], errors.UNSUPPORTED_PROTOCOL_VERSION)
    self.assertEqual(resp["error"]["data"]["requested"], "1999-01-01")
    self.assertIn(PROTOCOL_VERSION, resp["error"]["data"]["supported"])
\end{py}

\subsection{Statelessness}

Hard to test, and worth the effort, because a stateful optimisation is a
plausible thing for a colleague to add.

\begin{py}
def test_two_requests_share_no_state(self):
    """The second request must not benefit from the first having happened."""
    first = self.server.handle(request("tools/list", req_id=1))
    second = self.server.handle(request("tools/list", req_id=2))
    self.assertEqual(first["result"]["tools"], second["result"]["tools"])
    # And a request with no _meta still fails, even after a good one.
    third = self.server.handle(
        {"jsonrpc": "2.0", "id": 3, "method": "tools/list", "params": {}})
    self.assertIn("error", third)
\end{py}

\subsection{Error-code discipline}

Right codes, and equally important, no retired ones:

\begin{py}
def test_retired_codes_are_never_emitted(self):
    seen = set()
    for method, params in [("resources/read", {"uri": "x://y"}),
                           ("tools/call", {"name": "nope"}),
                           ("prompts/get", {"name": "nope"})]:
        resp = server.handle(request(method, params))
        seen.add(resp["error"]["code"])
    self.assertNotIn(-32002, seen)
    self.assertNotIn(-32042, seen)
\end{py}

\subsection{Caching hints}

All six operations, with \texttt{subTest} so one failure does not hide five
others.

\subsection{Header validation}

The transport-level contract, and each failure mode separately:

\begin{py}
def test_name_header_disagreeing_with_body_is_rejected(self):
    status, body = self.raw_post(
        {"jsonrpc": "2.0", "id": 1, "method": "tools/call",
         "params": {"name": "echo", "arguments": {"text": "hi"},
                    "_meta": meta()}},
        self.good_headers("tools/call", "some_other_tool"))
    self.assertEqual(status, 400)
    self.assertEqual(body["error"]["code"], errors.HEADER_MISMATCH)
\end{py}

Note \texttt{raw\_post}: a hand-built HTTP request. Testing header validation
through a client that constructs correct headers tests nothing.

\section{Integration testing with a stubbed model}

Now the whole system, deterministically.

The stub follows a script and has an LLM's cost profile but none of its
variance:

\begin{py}
class StubModel:
    """A planner that follows a script, with an LLM's cost profile.

    Deterministic given a seed, which is the only way the book's before/after
    numbers can mean anything.
    """
\end{py}

Which lets you assert things that would be impossible against a real model:

\begin{py}
def test_approval_flow_completes_in_two_round_trips(self):
    ...
    self.assertEqual(binding.client.stats.requests - before, 2)
    self.assertEqual(binding.client.stats.mrtr_rounds, 1)
\end{py}

\keyidea{Assert on \emph{round-trip counts}, not just on outcomes. A refactor
that makes a two-trip flow into a four-trip flow still returns the right answer,
and doubles the latency, and no correctness test will notice.}

Meridian tests the other direction too, because the cheap path being cheap is the
thing that regresses:

\begin{py}
def test_below_threshold_needs_no_round_trip(self):
    host.call_tool("risk.assess_account_risk",
                   {"accountId": "ACC-1000", "exposureUsd": 100_000})
    self.assertEqual(binding.client.stats.requests - before, 1)
\end{py}

\subsection{Testing performance properties}

You can test that parallel is faster than serial, deterministically, by injecting
latency:

\begin{py}
def test_parallel_fanout_beats_serial_on_round_trips(self):
    """Same work, same results, fewer wall-clock milliseconds."""
    host = wire_host()
    for binding in host.bindings.values():
        binding.client.transport.latency_ms = 8.0

    serial = AgentLoop(host, StubModel(plan=[calls, []],
                                       simulate_latency=False),
                       parallel_fanout=False).run("go")
    parallel = AgentLoop(host, StubModel(plan=[calls, []],
                                         simulate_latency=False),
                         parallel_fanout=True).run("go")

    self.assertLess(parallel.transport_ms, serial.transport_ms)
\end{py}

The assertion is a relation, not a threshold. \texttt{assertLess(x, 30)} fails on
a loaded CI runner; \texttt{assertLess(parallel, serial)} holds on any machine
fast or slow, because both sides move together.

\subsection{Testing the security controls}

Security tests must check both directions. A control that rejects everything
passes a naive test and breaks production.

\begin{py}
second = post({..., "requestState": forged}, 2)
self.assertIn("error", second, "a forged requestState was accepted")

# And the untampered state still works, so the check is not just
# rejecting everything.
third = post({..., "requestState": first["requestState"]}, 3)
self.assertIn("structuredContent", third["result"])
\end{py}

\section{Testing against real transports}

At least a few tests must use real sockets and real subprocesses, because that
is where a whole class of bug lives.

\begin{py}
def test_stdout_carries_only_mcp_messages(self):
    """One stray print in a dependency corrupts the stream for everyone."""
    transport = StdioClientTransport(
        [sys.executable, "-m", "meridian.servers.fraud"])
    try:
        client = Client(transport, server_label="fraud")
        for _ in range(5):
            client.call_tool("screen_account", {"accountId": "ACC-1001"})
        self.assertGreater(client.stats.requests, 4)
    finally:
        transport.close()
\end{py}

Subtle: the assertion is that five calls \emph{succeeded}. If anything had
written non-JSON to stdout, the reader would have discarded it and the calls
would have timed out. The test detects stream corruption without parsing
anything.

\begin{py}
def test_server_exits_when_stdin_closes(self):
    ...
    started = time.time()
    transport.close(timeout=6.0)
    self.assertLess(time.time() - started, 6.0,
                    "server had to be killed rather than exiting cleanly")
    self.assertIsNotNone(transport._proc.poll())
\end{py}

A server that ignores EOF gets SIGKILLed by every client it ever meets, and its
cleanup handlers never run. Nobody notices until a deploy leaves temp files
everywhere.

\section{Evals}

The top of the pyramid: does the thing actually work, with a real model.

Evals are not tests. They are a \emph{regression suite for behaviour}, and they
have properties tests do not: they are slow, they cost money, they are
non-deterministic, and their result is a score.

\subsection{What to measure}

\begin{table}[htbp]
\centering
\small
\begin{tabularx}{\textwidth}{@{}lX@{}}
\toprule
\thead{Metric} & \thead{What it catches} \\
\midrule
Task success rate     & The headline. Did it do the thing \\
Tool-selection accuracy & Catalogue bloat, ambiguous descriptions \\
Round trips per task  & Regressions in tool design or MRTR usage \\
Tokens per task       & Context bloat, catalogue growth \\
Cost per task         & All of the above, in the currency finance cares about \\
Recovery rate         & Whether tool errors are actually actionable \\
\bottomrule
\end{tabularx}
\caption{Six numbers. The last one is the one nobody measures and everybody
should.}
\label{tab:eval-metrics}
\end{table}

Recovery rate is worth explaining. Deliberately feed the agent a bad account id
and see whether it recovers. If your tool errors follow
Chapter~\ref{ch:tools}'s advice, the model reads the message, notices the format
hint, and retries correctly. If they say ``not found'', it gives up and asks the
user. That difference is invisible to every other metric.

\subsection{Building a suite}

\begin{py}
EVAL_CASES = [
    {
        "id": "single-account-risk",
        "goal": "What is the credit risk on ACC-1042?",
        "expect_tools": ["risk.assess_account_risk"],
        "expect_in_answer": ["55.4", "elevated"],
        "max_iterations": 3,
        "max_cost_usd": 0.02,
    },
    {
        "id": "portfolio-uses-batch-tool",
        "goal": "Rank ACC-1000 through ACC-1039 by risk.",
        "expect_tools": ["risk.rank_portfolio_risk"],
        "forbid_tools": ["risk.assess_account_risk"],   # the loop anti-pattern
        "max_iterations": 3,
    },
    {
        "id": "recovers-from-bad-id",
        "goal": "What is the risk on ACC-9999?",
        "expect_recovery": True,
        "max_iterations": 4,
    },
]
\end{py}

The \texttt{forbid\_tools} case is the useful one. It catches the model looping
over forty accounts one at a time, which is correct, slow, expensive, and exactly
what \texttt{rank\_portfolio\_risk} exists to prevent. Nothing else in your test
suite will notice.

\subsection{Scoring without flakiness}

Non-determinism makes naive assertions flaky. Three habits fix most of it:

\textbf{Assert on structure. ``The answer contains 55.4'' is stable; ``the answer says
the account is elevated risk'' is not.} ``The answer contains 55.4'' is stable.
``The answer says the account is elevated risk'' is not.

\textbf{Run N times, threshold the rate.} A case that passes 19 of 20 times is
passing. Requiring 20 of 20 is requiring determinism you were told you would not
get.

\textbf{Gate on movement, not on absolutes.} ``Success rate did not drop more
than 3 points against the last release'' is a gate you can actually keep green.
``Success rate is above 95\,\%'' is a gate somebody will eventually delete.

\section{Debugging}

Four failures cover most of what you will meet.

\subsection{``The client cannot connect''}

Almost always the era mismatch from Chapter~\ref{ch:architecture}. The client
sent \texttt{initialize} and got \texttt{-32601}, or sent modern
\texttt{\_meta} and got a legacy server's confusion.

\begin{sh}
printf '%s\n' \
 '{"jsonrpc":"2.0","id":1,"method":"server/discover","params":{"_meta":{"io.modelcontextprotocol/protocolVersion":"2026-07-28","io.modelcontextprotocol/clientCapabilities":{}}}}' \
 | python3 -m meridian.serve risk
\end{sh}

A \texttt{DiscoverResult} means modern. \texttt{-32601} means legacy. Anything
else means the server is broken in a more interesting way.

\subsection{``Everything times out on stdio''}

Something wrote to \texttt{stdout}. Check for a \texttt{print()}, a library
banner, or a warning. The tell is that the client hangs rather than erroring:
your line was consumed as a malformed message and no response was ever produced.

\begin{sh}
python3 -m meridian.serve risk < /dev/null 2>/dev/null | head -c 200
\end{sh}

Any output at all from a server that received no request is your culprit.

\subsection{``The model keeps picking the wrong tool''}

Three causes, in order of likelihood:

\begin{enumerate}[leftmargin=1.6em, itemsep=2pt]
\item The catalogue is too large. Measure \texttt{catalogue\_tokens()} and
compare against Chapter~\ref{ch:tools}.
\item Two descriptions do not distinguish themselves. Read them side by side and
ask which you would pick.
\item The right tool does not exist, and the model is approximating with what it
has. This is the most common one and the least often diagnosed.
\end{enumerate}

\subsection{``It works locally and not in production''}

Four candidates, all covered earlier and all worth a checklist:

\begin{itemize}[leftmargin=1.6em, itemsep=2pt]
\item Per-process MRTR signing secret behind a load balancer, so retries fail
about $1 - 1/N$ of the time.
\item In-memory task store behind a load balancer, same arithmetic.
\item A proxy buffering SSE because \texttt{X-Accel-Buffering: no} is missing.
\item A gateway rejecting the required \texttt{Mcp-*} headers because nobody told
it about them.
\end{itemize}

The shared symptom is \emph{intermittency proportional to replica count}, which
is a strong hint on its own.

\section{The MCP Inspector}

The official interactive tool, and the right first move when a server is
misbehaving and you do not know why. Connect, browse the catalogue, call tools by
hand, watch the raw messages.

What it gives you that a test does not: the ability to poke at the thing without
writing code first, and to see the exact bytes.

What it does not give you: a regression suite. Everything you learn in the
Inspector should end up as a contract test, or you will learn it again.

\begin{notebox}{Version skew in tooling}
Tooling lags the specification, sometimes by months. An inspector that expects an
\texttt{initialize} handshake will not connect to a strictly modern server, and
the failure looks like your server is broken.

This is another argument for the dual-era bridge: it makes your server debuggable
with today's tools while it stays correct for tomorrow's.
\end{notebox}

\section{CI}

What Meridian gates on, and roughly what I would gate on generally:

\begin{table}[htbp]
\centering
\small
\begin{tabularx}{\textwidth}{@{}lXl@{}}
\toprule
\thead{Gate} & \thead{Checks} & \thead{Runs on} \\
\midrule
Unit and contract & 191 tests, no model & every commit \\
Benchmarks        & Round trips and token counts did not regress & every commit \\
Prose lint        & Em dashes, slop, repetition, listing encodings & every commit \\
Reference check   & No \texttt{??} in the finished book & every commit \\
Evals             & Task success against the last release & every deploy \\
\bottomrule
\end{tabularx}
\caption{The middle row is unusual and I recommend it: assert that round-trip
counts and token costs have not gone up.}
\label{tab:ci-gates}
\end{table}

That benchmark gate deserves a word. A performance regression in an agent system
does not look like a slow function. It looks like one more round trip, or two
thousand more tokens per turn, and both are invisible to correctness tests and
obvious to a benchmark that runs on every commit.

\begin{sh}
make test      # 191 tests, about two seconds
make bench     # regenerates every number in the book
make lint      # prose rules plus cross-reference integrity
\end{sh}

\section{What to remember}

\begin{itemize}[leftmargin=1.6em, itemsep=3pt]
\item Four layers. Only the top needs a model, and only the top is slow.
\item An in-process transport gives fast tests and a benchmark control group
from one class.
\item Assert on the wire, not on convenience objects.
\item Assert on round-trip counts, in both directions.
\item Security tests must prove the control accepts valid input too.
\item Keep some tests on real sockets and real subprocesses. The stdout test
catches stream corruption without parsing anything.
\item Evals measure success rate, tool-selection accuracy, round trips, tokens,
cost, and recovery rate. Gate on movement, not absolutes.
\item Intermittency proportional to replica count means per-process state.
\item Everything you learn in the Inspector becomes a contract test.
\end{itemize}

Next: measuring the system you just learned to test.
